\section{完整模型的超图关联表示}
\label{sec:05_10}

第\ref{sec:05_08}节将任务、运动、观测、会合和能源条件合并为联合可行域，第\ref{sec:05_09}节进一步限定了信息可得性及状态接续条件。超图关联表示关注其中另一层关系：一条约束需要读取哪些决策变量，同一个变量又同时受到哪些约束限制。它为追踪模型的耦合位置提供结构表示，但并不替代约束的函数表达、变量域和物理单位。仅知道一条超边连接若干变量，尚不能判断这些变量应当满足等式、不等式、逻辑蕴含还是连续时间上的安全条件。

变量作为节点、约束作用域作为超边的构造，与约束满足问题中的依赖超图定义相联系\cite{Marx2011}。本节仅采用这一关联定义。巡检模型包含连续轨迹、整数决定、非线性函数和条件约束，不能据此直接移用离散约束满足问题的复杂性结论，也不由这种表示引入新的求解算法。

\subsection{连续函数依赖与有限标量依赖}

连续模型中的一个决策对象可能是一整段函数。例如，航次$\ell$的位置$\mathbf p_\ell(t)$是以$\mathrm m$计量的三维函数，物理电池对$p\in\mathcal B$的电量$E_p(t)$是以$\mathrm{Wh}$计量的标量函数。它们分别受到运动方程、观测质量、电量守恒等条件限制。此时可以说某个约束族依赖位置函数或电量函数，但不能将一个函数块直接当作一个有限数值变量，也不能将对全部时刻成立的约束族当作已经生成的一条标量约束。

令$\mathscr X$表示前述联合模型的决策空间，包含离散变量、事件时间及连续状态和控制函数。固定一种有限转录配置，记为$N$。这里$N$是配置标识，不是任务数；它规定有限事件槽、分段方式、函数参数化和已经选定的数值表达。该配置下的标量变量集合与取值向量分别记为
\begin{equation}
 \mathcal X_N=\{\chi_j:j\in\mathcal J_N^X\},
 \qquad
 \mathcal J_N^X=\{1,\ldots,n_X\},
 \qquad
 \mathbf x=(x_j)_{j\in\mathcal J_N^X}\in\mathcal D_N.
 \label{s0510:scalar_variables}
\end{equation}
$\chi_j$是变量身份，$x_j$是该变量的取值，索引$j$无量纲。集合$\mathcal D_N$规定各标量的基本取值域，以及表达求值所必需的定义域条件，不预先要求全部联合约束成立。二元分量取值于$\{0,1\}$，整数分量取值于相应整数集合，连续分量取值于规定的实数区间。位置参数、时间参数、电量参数分别采用$\mathrm m$、$\mathrm s$和$\mathrm{Wh}$；姿态与相机参数分别沿用第\ref{sec:04_04}节和第\ref{sec:04_05}节的表示及单位。因此，$\mathbf x$属于具有不同物理单位的乘积空间，不能给整个向量指定一个共同物理单位。

有限参数到连续对象的重建关系可写为
\begin{equation}
 \mathcal T_N:\mathcal D_N\longrightarrow\mathscr X,
 \qquad X=\mathcal T_N(\mathbf x).
 \label{s0510:reconstruction}
\end{equation}
$\mathcal T_N$恢复相应离散决定、事件时间以及所采用参数化下的轨迹和控制函数。该映射描述表达方式，并不保证其像覆盖原连续模型的全部可行解，也不保证每个$\mathbf x\in\mathcal D_N$重建出的对象均可行。对全部时间、视线参数或障碍物参数成立的条件，仍须保留原量词或由具有明确适用条件的数值检查处理，不能因为已经建立有限关联图就视为满足。

例如，一个三维轨迹参数块包含多个标量坐标；一个任务又可能关联分配、起止时间、轨迹、相机和能源等多类变量。区域观测任务数、变量块数量、有限标量变量数是三个不同层次的量。神经网络隐藏向量的维数属于算法内部表示，也不能替代这里的$n_X$。

\subsection{约束身份、作用域与关联矩阵}

令$\mathcal C_N$为所讨论有限表示中的约束索引集合。每个$c\in\mathcal C_N$具有唯一身份，其记录至少包括所属约束族、对象索引、适用条件、表达式和关系类型。对一个普通标量数值行，可统一写为
\begin{equation}
 \underline f_c\le F_c(\mathbf x;\theta_N)\le\overline f_c,
 \qquad c\in\mathcal C_N.
 \label{s0510:row_form}
\end{equation}
$\theta_N$是固定输入，包括几何、设备参数及当前允许使用的天气信息；它不属于决策变量集合。$\underline f_c$与$\overline f_c$为常数界，允许一端无界；等式对应两端相等。$F_c$及其有限界必须采用相同单位，不同行可以具有不同单位。逻辑蕴含若尚未转成这种数值行，则保留其逻辑关系，并将触发条件与被触发表达共同登记，不能只记录蕴含的后件。

区分函数意义上的依赖与固定表达中的记录依赖，有助于说明关联图能够支持什么结论。对定义在$\mathcal D_N$上的标量函数$F_c$，函数意义上的支持可以定义为
\begin{equation}
 \begin{split}
 S_c^{\mathrm{fun}}=\{\chi_j\in\mathcal X_N:\;&
 \exists\,\mathbf x,\mathbf x'\in\mathcal D_N,
 \quad x_h=x_h'\quad\forall h\ne j,\\
 &F_c(\mathbf x;\theta_N)\ne
 F_c(\mathbf x';\theta_N)\}.
 \end{split}
 \label{s0510:functional_support}
\end{equation}
这里比较的是在约束满足之前、函数有定义的候选取值。若先将比较域限制为满足$F_c=0$的解集，那么$F_c$在该集合上恒为零，却不能由此推断原等式不依赖任何变量。对于条件关系，也须考察完整关系的触发条件和表达，不能只在触发条件为假的候选点观察后件。

实际结构记录采用固定表达的作用域。记
\begin{equation}
 S_c=\operatorname{supp}_{\mathrm{rep}}(F_c)
 \subseteq\mathcal X_N,
 \qquad c\in\mathcal C_N,
 \label{s0510:recorded_support}
\end{equation}
其中下标$\mathrm{rep}$表示已经固定的表达及其变量叶收集规则。同一变量在一条表达中出现多次，只在$S_c$中登记一次；不同变量具有不同身份，即使其当前数值相同也不能合并。显式条件关系的作用域还包括决定触发与否的变量。式\eqref{s0510:recorded_support}记录表达可能读取的变量，不是某个候选点处非零偏导数对应的变量集合。

例如，令$q\in\mathbb R$为无量纲连续变量，$F(q)=q^2$。虽然在$q=0$处导数为零，改变$q$仍能改变函数值，故$q$不能从作用域中删除。二元触发变量更不能仅凭连续导数决定是否保留。另一方面，若表达中存在可以代数消去的重复项，直接遍历表达叶可能保留函数实际不再依赖的变量。因此，表达支持不自动等于经过任意代数消去后的最小函数支持。

已有表达记录实现逐行收集唯一标量变量叶，并在\texttt{Expr}项遍历中跳过零系数项；非线性表达则按所支持的运算递归读取变量叶。该过程没有给出任意代数化简后达到最小函数支持的证明。因此，本书对这种记录所得规模使用“固定表达下的结构量”，不将其解释为与表达无关的最小必要规模。

依赖超图及其关联矩阵定义为
\begin{align}
 \mathcal H_N^{\mathrm{dep}}
 &=\bigl(\mathcal X_N,(S_c)_{c\in\mathcal C_N}\bigr),
 \label{s0510:indexed_hypergraph}\\
 B_{jc}
 &=\mathbb I\{\chi_j\in S_c\},
 \qquad j\in\mathcal J_N^X,\quad c\in\mathcal C_N.
 \label{s0510:incidence_matrix}
\end{align}
$B_{jc}\in\{0,1\}$无量纲，$\mathbb I$为指标函数。这里使用带索引的超边族$(S_c)_{c\in\mathcal C_N}$，而不是将作用域收集成一个自动去重的集合。两条约束可以满足$S_c=S_{c'}$而$c\ne c'$；它们仍可能具有不同界、不同方向或不同适用条件，必须保留两列。即使两个表达完全相同，是否删除其中一条也属于显式模型处理，不能由超边集合去重暗中完成。

对于空作用域$S_c=\varnothing$，约束值只由固定输入决定。这样的行仍有约束身份：常数关系成立时它不限制决策，常数关系不成立时它可以直接否定相应模型或条件分支。空作用域不是遗漏约束的理由。

\subsection{结构量、目标依赖与七类联合关联}

在变量集合、约束集合和作用域规则均已固定后，定义
\begin{align}
 n_X&=|\mathcal X_N|,
 &m_C&=|\mathcal C_N|,
 \label{s0510:node_edge_counts}\\
 N_I&=\sum_{c\in\mathcal C_N}|S_c|
     =\sum_{j\in\mathcal J_N^X}\sum_{c\in\mathcal C_N}B_{jc},
 \label{s0510:incidence_count}\\
 r&=\max\bigl(\{|S_c|:c\in\mathcal C_N\}\cup\{0\}\bigr),
 \label{s0510:rank}\\
 d_j&=\sum_{c\in\mathcal C_N}B_{jc},
 &d_{\max}&=\max\bigl(\{d_j:j\in\mathcal J_N^X\}\cup\{0\}\bigr).
 \label{s0510:degrees}
\end{align}
以上均为无量纲非负整数。$n_X$是标量变量节点数，$m_C$是带唯一索引的约束行数，$N_I$是变量与约束的关联总数，$r$是一条约束所连接变量数的最大值，$d_j$是变量$\chi_j$关联的约束数，$d_{\max}$是最大变量度。$r$称为超图的秩，此处不表示几何半径。空约束集合给出$m_C=N_I=r=d_{\max}=0$；若变量集合也为空，则$n_X=0$。保留的变量即使没有连接任何约束，也仍计入$n_X$，其度为零。

目标依赖另行登记。令$J_{a,N}$为有限表示中第$a$个词典序目标表达，$a\in\{1,2,3\}$，则
\begin{equation}
 S_a^{\mathrm{obj}}
 =\operatorname{supp}_{\mathrm{rep}}(J_{a,N}),
 \qquad a\in\{1,2,3\}.
 \label{s0510:objective_support}
\end{equation}
第一项目标为无量纲未完成塔损失，等权时为未完成塔数；后两项分别以$\mathrm s$和$\mathrm{Wh}$计量，沿用第\ref{sec:05_08}节的词典序，不能因建立关联图而相加成未经定义的单一权重目标。目标作用域不默认并入$m_C$、$N_I$和变量度。若采用式\eqref{s0508:robust_epigraph}的上界变量，则目标表达可以只读取相应上界变量，而连接上界与各情形成本的关系作为约束行登记。此时关联路径改变，但目标的物理意义不变。一个约束度为零的变量仍可能进入目标，因而“孤立变量”只能在明确是否计入目标关联后解释。

表\ref{s0510:coupling_families}给出完整联合模型中需要保留的七类关联。表中列的是关联对象和作用，不将一个约束族压缩成一条超边；实际超边仍由每个标量约束索引逐条确定。

\begin{table}[htbp]
 \centering
 \small
 \caption{联合模型的七类变量关联及其约束作用}
 \label{s0510:coupling_families}
 \begin{tabular}{lll}\toprule
约束族 & 变量组 & 主要连接\\\midrule
整塔与任务流 & 完成、分配、访问弧 & 覆盖与次序\\
飞行与相机 & 位置、姿态、旋翼、云台 & 动力学与成像\\
道路与会合 & 道路、车辆状态、时刻 & 同时同地回收\\
时空冲突 & 轨迹、模式、操作时段 & 分离与容量\\
电池与充电 & 身份、SOC、极化、温度 & 状态与通道\\
天气与功率 & 轨迹、时刻、旋翼负载 & 共同情形\\
重规划与终态 & 观测、承诺、剩余服务 & 因果与接续\\\bottomrule
\end{tabular}
\end{table}

整塔与任务流关联可以直接从第\ref{sec:05_01}节和第\ref{sec:05_02}节读取。例如，整塔完成的下界行为
\begin{equation}
 z_g-1+\sum_{i\in\mathcal I_g}(1-A_i^k-a_i)\ge0,
 \qquad g\in\mathcal G.
 \label{s0510:tower_scope_example}
\end{equation}
$z_g,a_i\in\{0,1\}$分别表示计划整塔完成与本次任务完成决定，$A_i^k\in\{0,1\}$是边界时已验收的固定记录，全部无量纲。在保留这些决策变量的表达下，该行连接$z_g$与同塔全部$a_i$，固定记录$A_i^k$不作为变量节点。若某些变量已被常数替换，则替换后的作用域相应改变。任务分配行则连接同一任务在不同航次上的分配变量；访问流和时间衔接关系进一步将这些分配连接到服务顺序。只保留整塔一行，无法排除没有合法访问和回收路径的任务组合。

连续飞行与相机关联把几何可见性连接到可以执行的运动。第\ref{sec:05_03}节中的位置、速度和飞行内部加速度连接条件读取相邻轨迹段参数及段时长；曝光条件还读取姿态、云台和焦距。分段时钟中的时长影响时间导数，因此即使位置参数不变，改变时长也可能改变速度、模糊和能耗。固定目标样点是几何输入，不应被误计为决策节点。起落接触处仍采用既有单侧受力接口，关联表示不额外施加跨接触的全部高阶导数连续条件。

道路与会合关联连接车辆运动和飞行端点。会合位置关系中的飞机位置取自其到达时刻，车辆位置也必须在同一时刻求值；当到达时刻是变量时，它本身属于该关系的依赖。若采用停靠选择或道路分支变量，则相应触发变量也进入作用域。把车辆轨迹独立画出而不保留这些共同时间和位置的关联，不能保证飞机能够在所选载车上起降。

时空冲突关联同时读取几何状态与活动条件。两架飞机的分离约束只在相应飞行区间重叠时要求成立，但哪些区间重叠又由航次启用和起止时间决定。资源容量行可能同时连接多项操作的占用变量或时段端点。这里的关联含义是共同参与一个限制，并不等于这些动作在整个班次中永久互斥；同一资源在不相交时段可以顺序复用。

电池与充电关联必须保留物理身份。以第\ref{sec:05_05}节的电热转移为例，行内不仅有区间两端的SOC、极化和温度，还包括电流、充电输入、相应航次的电池分配、旋翼功率和时间边界。固定效率$\eta_p^{\mathrm{ch}}$是参数；输入功率以$\mathrm W$计量，时间积分以$\mathrm s$计量，除以$3600$后才与$\mathrm{Wh}$电量相加减。按总电量聚合而消去物理身份，会丢失同一电池不能同时飞行和充电、不同航次必须顺序接续的限制。关联图中的一条总能量行不能替代这些身份与模式行。

天气与能耗关联需要区分外生数据和决策自变量。给定情形$\xi\in\Xi_k$的天气场本身不是优化变量，但求值$\rho^\xi(\mathbf p_\ell^\xi(t),t)$仍可能依赖位置和时刻；由此得到的空速、密度和功率表达必须继续追踪其决策依赖。两个航次读取同一预报记录，并不自动产生资源冲突超边，因为预报数据不是被一个航次使用后就不能再被另一个航次使用的容量资源。不同情形的决定若受到非预见性关系限制，则通过实际的共同历史控制约束连接，而不是通过共享一个天气文件建立连接。

重规划与终态关联连接当前边界和未来动作。测得位置、电量和已验收状态在本次规划中作为固定输入，相应约束将未来轨迹或电量的初值绑定到这些输入；仍有效的承诺继续限制后续决定。返库、全部飞机和电池回到所属车辆，以及完成时刻以后不再补能，共同连接计划尾部。观测中途服务使用第\ref{sec:05_07}节的固定身份与虚拟源点关系，待验收图像使用接收时间门；这些实际方程的变量叶也必须进入关联，不能只给状态对象添加一个文字标签。

\subsection{条件分支、延迟生成与已注册结构}

一条条件约束在某个候选解处没有触发，不等于它已从模型中删除。设$b\in\{0,1\}$为无量纲触发变量，$G(\mathbf x)\le0$为具有一致单位的关系，则
\begin{equation}
 b=1\ \Longrightarrow\ G(\mathbf x)\le0
 \label{s0510:conditional_relation}
\end{equation}
的完整作用域包括$b$以及后件可能读取的变量。只有在$b$已经作为固定输入被消元，并且相应分支处理有明确定义时，才能据所选分支生成简化表示。若$b$仍是决策变量，不能因为当前候选取$b=0$就从完整候选域的依赖中删除后件。还须区分“被启用的模型行”与“在某个解处取等号的不等式”：后者是数值解上的紧约束概念，不能用来定义模型结构规模。

延迟生成则涉及另一层集合关系。对固定的有限候选表达域，令$\mathcal C_N^{\mathrm{gen}}$表示可生成的逻辑约束身份集合。计算阶段$h\in\mathbb Z_{\ge0}$已经生成的逻辑身份构成子集，但同一逻辑约束可能经历不同数值表达的替换。因此，再以$\mathcal C_{N,h}^{\mathrm{reg}}$记录截至阶段$h$累计保留的唯一注册行身份，并以$\varepsilon_c(h)\in\{0,1\}$表示该注册行是否仍启用。当前启用行集合为
\begin{equation}
 \mathcal C_{N,h}^{\mathrm{on}}
 =\{c\in\mathcal C_{N,h}^{\mathrm{reg}}:
                    \varepsilon_c(h)=1\}.
 \label{s0510:enabled_rows}
\end{equation}
$h$和$\varepsilon_c$均无量纲。可生成逻辑域、累计注册行和当前启用行不是同一个集合。未物化的选择分支尚不具有已记录的标量作用域；连续时间检查族还可能具有无限参数域，也不能仅由有限的$\mathcal C_N^{\mathrm{gen}}$替代。

已有表达带的结构统计使用保留的变量登记表、全部注册行以及各行的变量叶，并未按启用标记过滤。按其实际含义，式\eqref{s0510:node_edge_counts}--\eqref{s0510:degrees}用于这类输出时，应取
\begin{equation}
 \mathcal X_N=\mathcal X_{N,h}^{\mathrm{reg}},
 \qquad
 \mathcal C_N=\mathcal C_{N,h}^{\mathrm{reg}},
 \label{s0510:registered_convention}
\end{equation}
其中$\mathcal X_{N,h}^{\mathrm{reg}}$为该阶段仍在登记表中的标量变量身份集合。于是所报$m_C$、$N_I$、$r$、变量度及孤立变量数均描述累计保留的注册结构，可以包含已经停用的行。表达运算节点总数还包含常量和运算节点，与标量变量节点数$n_X$不同。

这一差别在能量积分表达更新时有实际含义：生成新的分段功率积分表达后，旧的能量连接行可以停用，再添加新的连接行。旧行的名称、变量关联和表达节点仍保留在记录中。求值和投影跳过停用行，但累计结构计数仍包含它们。因此，注册行数增长不一定表示当前有效约束数增加；某变量在注册结构中非孤立，也不保证它仍连接当前启用行。

若需要当前启用结构，应在同一变量登记口径下重新计算，例如
\begin{align}
 m_C^{\mathrm{on}}
 &=|\mathcal C_{N,h}^{\mathrm{on}}|,\\
 N_I^{\mathrm{on}}
 &=\sum_{c\in\mathcal C_{N,h}^{\mathrm{on}}}|S_c|,\\
 d_j^{\mathrm{on}}
 &=\sum_{c\in\mathcal C_{N,h}^{\mathrm{on}}}
       \mathbb I\{\chi_j\in S_c\}.
 \label{s0510:enabled_statistics}
\end{align}
这些量均无量纲，最大秩和最大度采用前述包含零的最大值约定。它们是启用行统计的定义，不能由现有累计输出直接冒充；现有材料未提供相应的启用行重计数。本节不填入未经核验的当前规模。运行结构统计须按累计注册、当前启用、未生成候选和连续检查族分别解释，不能把一个计数表理解为所有理论约束均已完全展开。

\subsection{表示改变与结构比较的条件}

即使可行决策含义不变，辅助变量也能改变关联统计。考虑一条标量关系
\begin{equation}
 F(\mathbf x)+G(\mathbf x)=0,
 \label{s0510:unsplit_expression}
\end{equation}
其中$F$和$G$在共同候选域上有定义，并具有相同物理单位。引入具有该单位的连续辅助变量$y$后，可将其写为
\begin{equation}
 y=F(\mathbf x),\qquad y+G(\mathbf x)=0.
 \label{s0510:auxiliary_expression}
\end{equation}
若$y$的域覆盖$F$在所讨论域上的全部取值，且没有额外限制，则消去$y$恢复式\eqref{s0510:unsplit_expression}；反过来，原关系的任一可行$\mathbf x$都可以令$y=F(\mathbf x)$而扩展为新表示的可行点。这个等价性来自明确的代入关系，不来自超图形状。

原行的作用域由$F$和$G$的记录支持合并得到，新表示则增加一个变量和一条约束，并形成分别连接$y$与两部分表达的超边。最大秩可能改变，关联总数和变量度也可能改变。如果给$y$施加过窄的界，上述投影等价性还会失效。因此，比较结构之前必须固定辅助变量、消元、表达化简、标量展开和约束分行规则；比较已有运行记录时还须固定物化阶段与停用行处理方式。

相同的变量块模型，采用不同轨迹段数、检查点、情形复制或积分表达，得到的$n_X$和$m_C$也可能不同。这些差异首先反映转录和记录口径，不能单独用于评价算法是否更好。模型完整性还需要将每个理论约束族对应到已生成关系、仍保留的连续检查或明确未闭合的接口；只统计关联数量无法发现物理条件是否被遗漏。

\subsection{与资源容量超图的区别及推论边界}

第\ref{sec:02_03}节的超图边选择模型以资源为顶点、候选任务为超边，变量$x_e\in\{0,1\}$表示是否选择候选任务，并由资源容量限制同时选择的组合。其保留实验口径为$n=|E|$个候选任务、$m=|V|$个资源。本节的$\mathcal H_N^{\mathrm{dep}}$则以标量决策为顶点、约束行为带索引超边，使用$n_X$和$m_C$避免混淆两种含义。

依赖超边不是可以任意选取或丢弃的任务。它记录某条必须满足的模型关系所读取的变量；条件约束的触发也由明确逻辑规定，不能交由一个边匹配算法随意删除。将关联矩阵转置，或者将共同出现在约束中的变量两两连边，都不会自动得到与原巡检模型等价的容量边选择问题。

若要建立这样的转换，至少需要给出原可行决策与转换后可行解之间的恢复关系，证明既不丢失所承诺域中的原可行决策，也不引入无法恢复为原可行决策的组合，并保持第\ref{sec:05_08}节三个目标的取值或词典序比较。轨迹、会合和电池接续等连续可行性必须包含在这种证明内。仅凭候选之间共用资源，无法证明一组被选中的候选具有联合可执行的连续见证。

现有纯组合实验全部采用单位资源容量，其二元可行域可以按第\ref{sec:02_04}节表示为两两冲突约束。因此，高秩本身不能证明不可约的高阶依赖。对本节的联合模型，一条约束涉及许多变量也只表明当前表达的作用域较大；是否存在保持域和目标的分解，以及这种分解是否有利于求解，需要另外分析。结构量不包含函数曲率、数值尺度、松弛强度或验收开销，不能单独推出速度或解质量优势。

超图关联表示在这里承担的是保留约束身份、追踪变量联系和解释结构统计的作用。对不同求解方法的比较，仍须固定几何输入、任务与道路域、信息截止、物理约束、初末状态和目标，并结合实际求解与验收证据。有限结构已经构建、候选通过局部或整体数值检查、取得完整覆盖、证明最优，是不同结论，均不能由超图的规模或连通性相互替代。
